\subsection{The ring of Witt vectors with coefficients in a ring of characteristic p}

PROPOSITION 5. --- Let A be a ring of characteristic p (A, V, p. 2). For any elements a and b of W(A), and positive integers m and n, if \(\boldsymbol{a} = (a_{n})_{n \in \mathbb{N}}\) ,

\begin{enumerate}
\def\labelenumi{(\arabic{enumi})}
\setcounter{enumi}{50}
\tightlist
\item
\end{enumerate}

\[
\mathrm{F} (\boldsymbol {a}) = (a _ {n} ^ {p}) _ {n \in \mathbf {N}}\tag{52}
\]

\[
p. \boldsymbol {a} = \operatorname{VF} (\boldsymbol {a}) = \operatorname{FV} (\boldsymbol {a}) = (0, a _ {0} ^ {p}, a _ {1} ^ {p}, \dots)\tag{53}
\]

\[
\mathrm{V} ^ {m} (a) \times \mathrm{V} ^ {n} (b) = \mathrm{V} ^ {m + n} (\mathrm{F} ^ {n} (a) \times \mathrm{F} ^ {m} (b)).
\]

Formula (51) follows from Example 4 of No. 3. From this we immediately deduce the equality

\[
\operatorname{VF} (\boldsymbol {a}) = \operatorname{FV} (\boldsymbol {a}) = (0, a _ {0} ^ {p}, a _ {1} ^ {p}, \dots),
\]

and the equality \(p\) . \(a = \text{FV}(a)\) was proved (no. 5, prop. 3), whence (52).

Let us prove (53). By formula (37) (where we substitute \(V^n(\boldsymbol{b})\) for \(\boldsymbol{b}\) ), we have

\[
\mathrm{V} ^ {m} (\boldsymbol {a}) \times \mathrm{V} ^ {n} (\boldsymbol {b}) = \mathrm{V} ^ {m} \left(\boldsymbol {a} \times \mathrm{F} ^ {m} \left(\mathrm{V} ^ {n} (\boldsymbol {b})\right)\right).\tag{54}
\]

From formula (37), we also deduce

\[
\mathrm{V} ^ {n} \left(\mathrm{F} ^ {m} (\boldsymbol {b})\right) \times \boldsymbol {a} = \mathrm{V} ^ {n} \left(\mathrm{F} ^ {m} (\boldsymbol {b}) \times \mathrm{F} ^ {n} (\boldsymbol {a})\right).\tag{55}
\]

Formula (53) then follows from (54) and (55) and the relation \(F^{m} \circ V^{n} = V^{n} \circ F^{m}\) , itself a consequence of (51).

COROLLARY. --- If m and n are two positive integers, we have

\[
\mathrm{V} _ {m} (\mathrm{A}) \times \mathrm{V} _ {n} (\mathrm{A}) \subset \mathrm{V} _ {m + n} (\mathrm{A}).
\]

This follows from formula (53), since \(V_{m}(A)\) is the image of \(V^{m}:W(A)\to W(A)\) .

PROPOSITION 6.---Let A be a ring.

\begin{enumerate}
\def\labelenumi{\alph{enumi})}
\item
  For every integer \(k \geqslant 1\) , one has \((\mathrm{V}_1(\mathrm{A}))^k = p^{k-1} \cdot \mathrm{V}_1(\mathrm{A})\) .
\item
  Suppose that A is a ring of characteristic p. On the ring W(A), the V₁(A)-adic topology and the p-adic topology coincide, and they are finer than the product topology ℰ (cf. no. 6). The ring W(A) is separated and complete for the p-adic topology.
\end{enumerate}

Let us prove \(a\) ) by induction on \(k\) . The case \(k = 1\) is evident. Suppose \(k \geqslant 2\) . By the induction hypothesis, we have \(\mathrm{V}_{1}(\mathrm{A})^{k-1} = p^{k-2} \cdot \mathrm{V}_{1}(\mathrm{A})\) and consequently \(\mathrm{V}_{1}(\mathrm{A})^{k} = p^{k-2} \cdot (\mathrm{V}_{1}(\mathrm{A}))^{2}\) . But it follows from prop. 3, \(d\) ), formula (31), of no. 5 that one has \((\mathrm{V}_{1}(\mathrm{A}))^{2} = p \cdot \mathrm{V}_{1}(\mathrm{A})\) , whence \(a\) ).

Suppose now that A has characteristic p. Since one has

\[
p. \mathrm{W} (\mathrm{A}) = \mathrm{VF} (\mathrm{W} (\mathrm{A})) \subset \mathrm{V} _ {1} (\mathrm{A}) \quad (\text { formula (52) }),
\]

one deduces from \(a\) ) the inclusions \(p^{k} \cdot \mathbf{W}(\mathbf{A}) \subset (\mathbf{V}_{1}(\mathbf{A}))^{k} \subset p^{k-1} \cdot \mathbf{W}(\mathbf{A})\) , and from the corollary to prop. 5 the inclusion \((\mathbf{V}_{1}(\mathbf{A}))^{k} \subset \mathbf{V}_{k}(\mathbf{A})\) , for every integer \(k \geqslant 1\) . The first assertion of \(b\) ) follows from this.

Let \(k\) be an integer \(\geqslant 1\) . By formula (52), the ideal \(p^{k}\) . W(A) of W(A) is the set of elements \(\boldsymbol{a} = (a_{n})_{n \in \mathbf{N}}\) of W(A) such that one has \(a_{n} = 0\) for \(n < k\) and \(a_{n} \in A^{p^{k}}\) for \(n \geqslant k\) . It is therefore closed for the topology \(\mathcal{G}\) . Since W(A) is separated and complete for the topology \(\mathcal{G}\) and the ideals \(p^{k}\) . W(A) of W(A), for \(k \geqslant 1\) , form a neighborhood base of 0 in W(A) for the \(p\)-adic topology, the ring W(A) is separated and complete for the \(p\)-adic topology (TG, III, p. 26, cor. 1 to prop. 10).

PROPOSITION 7. --- Let A be a perfect ring of characteristic p.

\begin{enumerate}
\def\labelenumi{\alph{enumi})}
\item
  For every element \(\boldsymbol{a} = (a_{n})_{n \in \mathbb{N}}\) of W(A), the series with general term \(p^{n}\tau(a_{n}^{p^{-n}})\) is convergent in W(A), with sum a.
\item
  On \(\mathbf{W}(\mathbf{A})\) , the \(\mathbf{V}_1(\mathbf{A})\)-adic topology, the \(p\)-adic topology and the topology \(\mathcal{C}\) coincide. More precisely, we have \(\mathbf{V}_n(\mathbf{A}) = p^n\) . \(\mathbf{W}(\mathbf{A}) = (\mathbf{V}_1(\mathbf{A}))^n\) for every integer \(n \geqslant 0\) . In particular \(\Phi_0\) defines an isomorphism from \(\mathbf{W}(\mathbf{A})/p\) . \(\mathbf{W}(\mathbf{A})\) onto \(\mathbf{A}\) .
\end{enumerate}

By definition (A, V, p. 5), the map \(a \mapsto a^{p}\) is an automorphism of the ring A. By prop. 5, F is therefore an automorphism of the ring W(A), and one has, for every \(n \in \mathbf{N}\) ,

\[
p ^ {n}. \mathrm{W} (\mathrm{A}) = \mathrm{V} ^ {n} \mathrm{F} ^ {n} (\mathrm{W} (\mathrm{A})) = \mathrm{V} ^ {n} (\mathrm{W} (\mathrm{A})) = \mathrm{V} _ {n} (\mathrm{A}).
\]

In particular, we have \((\mathbf{V}_1(\mathbf{A}))^n = (p.\mathbf{W}(\mathbf{A}))^n = p^n.\mathbf{W}(\mathbf{A})\) . The assertion \(b)\) follows from that. By prop. 5, we have

\[
p ^ {n} \cdot \tau (a _ {n} ^ {p ^ {- n}}) = \mathrm{V} ^ {n} \mathrm{F} ^ {n} \tau (a _ {n} ^ {p ^ {- n}}) = \mathrm{V} ^ {n} \tau (a _ {n}),
\]

and assertion a) follows from prop. 4 of no. 6.

PROPOSITION 8. — Let A be a field of characteristic p. The ring W(A) is a local integral domain, separated and complete, with maximal ideal V₁(A) and residue field isomorphic to A. If the field A is perfect, the ring W(A) is a discrete valuation ring, and its maximal ideal is p.W(A).

The homomorphism \(\Phi_{0}\) defines an isomorphism from W(A)/V \(_{1}\) (A) onto A (no. 7, example 1). The ideal V \(_{1}\) (A) of W(A) is therefore maximal. Since the ring W(A) is separated and complete for the topology V \(_{1}\) (A)-adic (prop. 6, b)), it is a local ring, with maximal ideal V \(_{1}\) (A) (III, § 2, no. 13, prop. 19).

Let \(\pmb{a}\) and \(\pmb{b}\) be two nonzero elements of \(\mathbf{W}(\mathbf{A})\) . There exist integers \(m \geqslant 0\) and \(n \geqslant 0\) , and elements \(\pmb{a}' = (a_n')_{n \in \mathbb{N}}\) and \(\pmb{b}' = (b_n')_{n \in \mathbb{N}}\) of \(\mathbf{W}(\mathbf{A})\) such that \(\pmb{a} = \mathbf{V}^m(\pmb{a}')\) , \(\pmb{b} = \mathbf{V}^n(\pmb{b}')\) and such that the elements \(a_0'\) and \(b_0'\) of \(\mathbf{A}\) are nonzero. Then the component of index \(m + n\) of \(\pmb{a} \times \pmb{b}\) is equal to the component of index 0 of \(\mathbf{F}^n(\pmb{a}') \times \mathbf{F}^m(\pmb{b}')\) (formula (53)), that is, to \(a_0'^{p^n} b_0'^{p^m}\) (formula (51) and no. 3, example 2). Consequently \(\pmb{a} \times \pmb{b}\) is nonzero and \(\mathbf{W}(\mathbf{A})\) is integral.

If the field A is perfect, the maximal ideal \(V_{1}(A)\) of \(W(A)\) is equal to p.\(W(A)\)

(prop. 7, b)) and consequently W(A) is a discrete valuation ring (VI, § 3, no. 6, prop. 9, c)).

Remarks. --- 1) Let A be a field of characteristic p. One can show that the ring W(A) is Noetherian if and only if A is perfect (p. 43, exerc. 9).

\begin{enumerate}
\def\labelenumi{\arabic{enumi})}
\setcounter{enumi}{1}
\tightlist
\item
  Let A be a ring of characteristic p. By prop. 5, we have the formulas
\end{enumerate}

\[
\mathrm{F} _ {m} ^ {n} [ a _ {0}, \dots , a _ {n + m - 1} ] = [ a _ {0} ^ {p ^ {n}}, \dots , a _ {m - 1} ^ {p ^ {n}} ]
\]

\[
p ^ {n} \cdot [ a _ {0}, \dots , a _ {n + m - 1} ] = [ \underbrace {0 , \dots , 0} _ {n \text {   times }}, a _ {0} ^ {p ^ {n}}, \dots , a _ {m - 1} ^ {p ^ {n}} ]
\]

for every Witt vector \([a_0, \dots, a_{n+m-1}]\) of length \(n + m\) .

In fact, the map \(F:W(A)\to W(A)\) allows one, by passing to quotients by \(V_{m}(A)\) , to define a map \(\overline{F}_{m}:W_{m}(A)\to W_{m}(A)\) . We have the formula

\[
\overline {{{{\mathrm{F}}}}} _ {m} \left[ a _ {0}, \dots , a _ {m - 1} \right] = \left[ a _ {0} ^ {p}, \dots , a _ {m - 1} ^ {p} \right].
\]

The maps \(V_{m}^{1} \circ \overline{F}_{m}\) and \(\overline{F}_{m+1} \circ V_{m}^{1}\) from \(W_{m}(A)\) to \(W_{m+1}(A)\) are equal and are deduced, by passing to the quotient, from multiplication by p in \(W_{m+1}(A)\) .

